\section{Conclusion}

Educational AI systems are often discussed as if retrieval quality alone determines usefulness. This paper argues for a more specific framing. The important design question is not simply whether a system retrieves documents, but whether a fixed retrieve-then-generate pipeline remains sufficient once the task demands pedagogical sequencing, critique, learner-state diagnosis, and adaptive support.

The article makes three contributions. It proposes a task taxonomy that separates knowledge-grounding need from pedagogical coordination complexity. It develops an evaluation framework that compares classical RAG, agentic RAG, and non-RAG multi-agent systems under shared tasks, budgets, and benchmark regimes while keeping groundedness, pedagogical quality, adaptivity, robustness, and cost analytically separate. It also adds repository-grounded evidence from EduBench and TutorEval snapshots showing that the preferred architecture changes with task regime: coordination-heavy conditions favor multi-agent or hybrid orchestration, while long-context evidence-grounded tutoring still benefits from conditional retrieval.

The resulting conclusion is intentionally disciplined. Multi-agent systems should not be presented as a blanket replacement for retrieval in education. Instead, they are best understood as a stronger orchestration option when pedagogical coordination becomes the dominant systems problem and when retrieval is more useful as a selectively invoked capability than as the sole top-level design pattern. A complete empirical version of the paper should extend the current snapshots to full benchmark runs, add architecture-blind human evaluation, and report confidence intervals and deployment costs. Even in its current form, however, the manuscript offers a clearer way to ask the comparison question and a more defensible basis for deciding when educational assistance requires more than classical RAG.
